\subsection{相位塔几何接口：主 \texorpdfstring{$U(1)$}{U(1)} 束、面积权重与精度放大律}\label{subsec:unit-disk-phase-tower}
本附录的“相位紧化门”可被置于一个更一般的几何接口中：每一次“提升”在结构上表现为一个主 \(U(1)\) 束（相位纤维），而“近零加权”则由均匀测度的 Jacobian 拉回自动产生。
本小节仅给出与本文门语言可直接拼接的最小形式。

\begin{remark}[标准链：射影化 \(\to\) Hopf 主 \texorpdfstring{$U(1)$}{U(1)} 束 \(\to\) 立体投影体积元拉回]\label{rem:phase-tower-standard-chain}
射影化与 Hopf 纤维化把“提升一次多出一个全局相位纤维”严格实现为主 \(U(1)\) 束；立体投影下体积元拉回给出 \((1+|x|^{2})^{-n}\) 型权重衰减。以上均为标准几何事实，见 \cite{Husemoller1994FiberBundles}。
\end{remark}

\paragraph{相位塔：主 \texorpdfstring{$U(1)$}{U(1)} 束与 Hopf 纤维化}
Hopf 纤维化给出标准例子
\[
S^{1}\hookrightarrow S^{3}\twoheadrightarrow S^{2},
\]
并对任意 \(n\ge 1\) 给出
\[
S^{1}\hookrightarrow S^{2n+1}\twoheadrightarrow \mathbb{CP}^{n}.
\]
该形式表达“提升一次，多出一个全局相位纤维”的结构原语（参见 \cite{Husemoller1994FiberBundles}）。
在本文语境中，\(\upsilon\) 与 \(\iota\) 所提供的相位门可以被视为把无界坐标压到一个紧相位纤维上的最小实例。

\begin{proposition}[射影化与主 \texorpdfstring{$U(1)$}{U(1)} 束（相位塔原语）]\label{prop:phase-tower-principal-u1}
对任意 \(n\ge 1\)，令
\(
S^{2n+1}=\{z\in\CC^{n+1}:\|z\|=1\}
\)，并定义投影
\[
p:S^{2n+1}\to \mathbb{CP}^{n},\qquad p(z):=[z]
\]
（\([z]\) 表示由 \(z\) 张成的复直线）。
则 \(U(1)\) 通过 \(e^{\mathrm{i}\phi}\cdot z:=e^{\mathrm{i}\phi}z\) 自由作用在 \(S^{2n+1}\) 上，且 \(p\) 是商映射；因此 \(p\) 构成一个主 \(U(1)\) 束。
当 \(n=1\) 时，\(p:S^{3}\twoheadrightarrow\mathbb{CP}^{1}\cong S^{2}\) 即 Hopf 纤维化。
\end{proposition}
\begin{proof}
相位群 \(U(1)\) 的作用保持范数，且 \(e^{\mathrm{i}\phi}z=z\Rightarrow e^{\mathrm{i}\phi}=1\)（自由性）。
两点 \(z,z'\in S^{2n+1}\) 具有相同射影像当且仅当 \(z'=\lambda z\)（\(\lambda\in\CC^\times\)）；由 \(\|z\|=\|z'\|=1\) 得 \(|\lambda|=1\)，故 \(\lambda\in U(1)\)。
因此 \(p\) 的纤维恰为 \(U(1)\)-轨道，\(p\) 为主 \(U(1)\) 束的标准模型（参见 \cite{Husemoller1994FiberBundles}）。证毕。
\end{proof}

\paragraph{局部相位与全局相位：射影化与 \texorpdfstring{$\det:U(n)\to U(1)$}{det}}
将多轴信息编码为复向量 \(z\in\CC^{n+1}\) 时，可对每个坐标写极坐标分解 \(z_k=r_k e^{\mathrm{i}\theta_k}\)，从而每一轴均携带一个\textbf{局部相位} \(\theta_k\)。
与此同时存在一个共同的\textbf{全局相位}作用 \(z\mapsto e^{\mathrm{i}\phi}z\)（\(\phi\in\RR\)），其商空间即为复射影空间：
\[
\mathbb{CP}^{n}\ \cong\ (\CC^{n+1}\setminus\{0\})/U(1).
\]
因此，“合并后只有一个角度”的硬化版本可以表述为：全局相位是主 \(U(1)\) 束的纤维坐标，而相对结构落在基空间 \(\mathbb{CP}^{n}\) 上。
在矩阵层面，任意 \(U\in U(n)\) 的行列式 \(\det(U)\in U(1)\) 给出一个唯一的“总体相位读出”，且其核为 \(SU(n)\)；因此 \(\det:U(n)\to U(1)\) 提供了一个代数化的全局相位抽取口径。

\paragraph{面积权重与 Jacobian：近零加权的结构来源}
在一维情形中，\(\upsilon(x)=\frac{1}{\pi}\arctan(x)\) 的导数给出自然权重 \(\dd\mu(x)=\frac{1}{\pi(1+x^2)}\dd x\)（见 \eqref{eq:unit-circle-phase-derivative}--\eqref{eq:unit-circle-phase-cauchy-measure}）。
该权重并非外加规则，而是紧化换元的 Jacobian。
在二维/更高维中，类似的现象由“均匀面积元/体积元在投影坐标下的密度因子”给出。
例如在 Riemann 球面 \(S^{2}\) 的立体投影坐标 \(z=X+\mathrm{i}Y\in\CC\) 下，圆度量满足
\[
ds^{2}=\frac{4}{(1+|z|^{2})^{2}}\,(\dd X^{2}+\dd Y^{2})
=\frac{4}{(1+|z|^{2})^{2}}\,|\,\dd z\,|^{2},
\]
从而球面均匀面积元的拉回密度为
\[
\boxed{\ \dd A_{S^{2}}=\frac{4}{(1+|z|^{2})^{2}}\,\dd A_{\RR^{2}}\ }.
\]
因此在平面坐标中，\(|z|\) 越大对应的球面面积越小；“近零加权”由 Jacobian 自动产生，而非额外规定。

\begin{proposition}[立体投影的体积元拉回（一般维度）]\label{prop:stereographic-volume-form}
\leanverified{paper\_stereographic\_volume\_form}
令 \(\sigma:S^{n}\setminus\{N\}\to\RR^{n}\) 为从北极 \(N\) 出发的立体投影，并以 \(x\in\RR^{n}\) 为像空间坐标。
令 \(\sigma^{-1}:\RR^{n}\to S^{n}\setminus\{N\}\) 为其逆映射。
则由 \(\sigma^{-1}\) 拉回的球面度量满足共形缩放
\[
(\sigma^{-1})^{\ast}g_{S^{n}}=\left(\frac{2}{1+\|x\|^{2}}\right)^{2}g_{\mathrm{Euc}},
\]
从而球面体积形式的拉回满足
\begin{equation}\label{eq:stereographic-volume-general}
\boxed{\;
(\sigma^{-1})^{\ast}\dd V_{S^{n}}
 =\left(\frac{2}{1+\|x\|^{2}}\right)^{n}\dd V_{\RR^{n}}
=\frac{2^{n}}{(1+\|x\|^{2})^{n}}\,\dd x_{1}\cdots \dd x_{n}\; }.
\end{equation}
特别地，\(n=1\) 时该密度与 \(\dd\mu(x)=\frac{1}{\pi(1+x^{2})}\dd x\) 仅差一个归一化常数，而 \(n=2\) 时退化为上式的 \(\frac{4}{(1+\|x\|^{2})^{2}}\)。
\end{proposition}
\begin{proof}
立体投影给出球面在 \(\RR^{n}\) 上的标准共形坐标；共形因子为 \(\Omega(x)=\frac{2}{1+\|x\|^{2}}\)，即 \((\sigma^{-1})^{\ast}g_{S^{n}}=\Omega(x)^{2}g_{\mathrm{Euc}}\)。
体积形式在共形缩放下按 \(\Omega^{n}\) 缩放，因此
\(
(\sigma^{-1})^{\ast}\dd V_{S^{n}}=\Omega(x)^{n}\dd V_{\RR^{n}}
\)，得到 \eqref{eq:stereographic-volume-general}。证毕。
\end{proof}

\paragraph{精度放大律：降维紧化把动态范围压入末位精度}
由 \(\upsilon^{-1}(u)=\tan(\pi u)\)（式 \eqref{eq:unit-circle-phase-inverse}）得误差放大
\begin{equation}\label{eq:app-precision-amplification}
\boxed{\ \frac{\dd x}{\dd u}=\pi(1+x^2)\ }.
\end{equation}
因此当 \(|x|\) 大时，微小的相位误差 \(\Delta u\) 将被放大为 \(\Delta x\approx \pi x^2\Delta u\) 的巨大读出误差；换言之，降维/紧化把“远离零点的大数信息”强制搬运到 \(u\) 的末位精度中。
粗略地，若 \(u\) 的量化精度为 \(\Delta u\)，则可分辨到的最大量级满足
\[
|x|_{\max}\approx \cot(\pi\Delta u)\approx \frac{1}{\pi\Delta u}\qquad(\Delta u\to 0).
\]
这为“信息被压到精度上”的门语言给出一条可核验的解析尺度律。
